\section{Introduction}
\label{sec:intro}

Recent advances in Vision-Language-Models~(VLM) brought LLM-type scaling and reasoning
into the domain of robotics. As VLMs are used in both next-action prediction
(Vision-Language-Action models) and higher-level task planning and supervision, they
require state understanding of the robot and the objects it interacts with, given the
context of the task. To bridge the gap to real-world deployment, the models have to be 
self-sufficient, thus generalizing to unforeseen
environments, tasks, and edge cases, relieving the user from needing to finetune the
model excessively. Thus it is necessary for these models to reliably track progress during a
task, identify which sub-task the robot is currently attempting, and understand
success/failure of robotic actions. Furthermore it has to be time and resource efficient
based on the task and an overall pipeline understanding: e.g.\ for a household robot it
might suffice to empty the dishwasher under 20 minutes, while an industrial production
pipeline might have only several seconds to determine success or failure of a task. We
investigate the zero-shot capabilities of different model families and sizes for failure
mode understanding and multi-view consistency, to determine under which task-specific and
scene conditions the models' reasoning collapses.

%TODO
We find that current small- and mid-sized VLMs fail to deliver this level of reliability:
zero-shot accuracy stays close to random across model families, and explicitly prompting
models to reason step by step before answering makes performance worse, not better. While
such models exhibit robust performance across object detection, positional understanding, they struggle with object states and complex cross-object relationships.  robot supervision poses a
narrower and harder questions: Robots require supervision beyond raw
perception: detecting failures, tracking progress, and reasoning about the current action
phase are prerequisites for any higher-level planning or human-in-the-loop monitoring
built on top of a perception stack. The open question we address is at which task- and scene-complexity, models fail. 

We study this gap along two concrete capability questions. First, failure-mode and state
reasoning: given multiple views of the scene, can a VLM judge task progress,
identify the current action phase, and determine whether action-phase execution succeeded or failed?
Second, multi-view consistency: do answers extracted from different camera angles of the
 moment agree with each other? Inconsistency across views would
undermine true scene understanding and entails intensive calibration phase.

To answer these questions, we build a new benchmark on
Robo2VLM-1~\cite{chen2025robo2vlm} (DROID~\cite{khazatsky2024droid} and OXE
trajectories), covering 47 scenes and 6224 questions across four question types: four
failure-mode sub-tasks and one multi-view consistency task. We evaluate seven VLMs across
three families zero-shot, ablate chain-of-thought (CoT) prompting, analyze accuracy as a
function of task- and scene-difficulty, and run a LoRA fine-tuning case study on Qwen3-4B with a
decomposition of where the performance gain actually come from.
\begin{itemize}
  \item New benchmark: 6224 questions / 47 scenes on Robo2VLM-1 (DROID), covering
        failure-mode detection (4 sub-tasks) and multi-view consistency.
  \item Self-designed question types (action phase, phase success, progress, task success,
        multi-view consistency), including adversarial ``Cannot be determined'' (CnBD) variant.
  \item Zero-shot evaluation of 7 VLMs across 3 families (Gemma3, Qwen3, InternVL3) with
        easy/hard scene-difficulty breakdown and CoT ablation.
  \item LoRA fine-tuning case study (Qwen3-4B): +10.9pp overall, with detailed gain
        decomposition revealing two coarse heuristics rather than improved phase reasoning.
\end{itemize}

% --- Figure: method overview ---
\begin{figure}[t]
  \centering
  \includegraphics[width=\linewidth]{../poster_3_D/method_overview.pdf}
  \caption{Benchmark pipeline: trajectories from Robo2VLM-1 (DROID/OXE) are turned into
  multiple-choice failure-mode and multi-view consistency questions, which are posed to
  seven zero-shot VLMs, a CoT variant, and a LoRA fine-tuned model, then scored and
  decomposed by question type and answer class.}
  \label{fig:method_overview}
\end{figure}
